\documentclass{article}
% Encoding and hyphenation
\usepackage[T1]{fontenc}
\usepackage[utf8]{inputenc}
\usepackage[spanish,english,shorthands=off]{babel}
\usepackage{graphicx}
\usepackage[table]{xcolor}
\usepackage{array}
\usepackage{ragged2e}
\usepackage{enumitem}
\usepackage{hyperref}
\usepackage{xparse}
\graphicspath{{src/img/}{img/}}
\makeatletter
\def\input@path{{}{src/}}
\makeatother
% Descripción de un requerimiento
% Ejemplo de uso:
% \req[
%   identificador=RNF-RE-04,
%   titulo={Red de distribución de contenido},
%   descripcion={Los ficheros estáticos...},
%   prioridad=Could Have
% ]
\ExplSyntaxOn
\keys_define:nn { req }
  {
    identificador .tl_set:N = \l__req_id_tl,
    titulo .tl_set:N = \l__req_title_tl,
    descripcion .tl_set:N = \l__req_description_tl,
    prioridad .tl_set:N = \l__req_priority_tl
  }
\NewDocumentCommand{\req}{ O{} }
  {
    \keys_set:nn { req } { #1 }
    {\small
    \par\vspace{0.8\baselineskip}\noindent\phantomsection\hypertarget{req:\l__req_id_tl}{}\label{req:\l__req_id_tl}%
    \textbf{\l__req_title_tl}\par\vspace{0.3\baselineskip}%
    \noindent\begin{tabular}{@{}>{\bfseries\RaggedRight\arraybackslash}p{4.0cm}p{\dimexpr\linewidth-4.0cm-2\tabcolsep\relax}@{}}
    \arrayrulecolor{lightgray}
    Identificador: & \texttt{\l__req_id_tl} \\
    \hline
    Descripción: & {\l__req_description_tl} \\
    \hline
    Prioridad~(MoSCoW): & {\l__req_priority_tl} \\
    \end{tabular}\par
    }% end small
    \vspace{0.6\baselineskip}%
  }
\ExplSyntaxOff

% Referencia a un requerimiento por su identificador con un enlace clicable
% Ejemplo de uso: \reqref{RNF-SE-01}
\newcommand{\reqref}[1]{\hyperref[req:#1]{\texttt{#1}}}

% Historia de usuario: identificador, título, rol (como), objetivo (quiero), puntos, criterios, tareas
% Ejemplo de uso:
% \us[
%   identificador=US-XX-YY,
%   titulo={Mi historia},
%   como={Rol},
%   quiero={Objetivo},
%   puntos={3},
%   criterios={...},
%   tareas={...}
% ]
\ExplSyntaxOn
\keys_define:nn { us }
  {
    identificador .tl_set:N = \l__us_id_tl,
    titulo .tl_set:N = \l__us_title_tl,
    como .tl_set:N = \l__us_role_tl,
    quiero .tl_set:N = \l__us_goal_tl,
    requerimientos .tl_set:N = \l__us_requirements_tl,
    puntos .tl_set:N = \l__us_points_tl,
    criterios .tl_set:N = \l__us_acceptance_tl,
    tareas .tl_set:N = \l__us_tasks_tl
  }
\NewDocumentCommand{\us}{ O{} }
  {
    \keys_set:nn { us } { #1 }
    \par\vspace{0.8\baselineskip}\noindent\phantomsection\hypertarget{us:\l__us_id_tl}{}\label{us:\l__us_id_tl}%
    \textbf{\l__us_title_tl}\par\vspace{0.3\baselineskip}%
    \noindent\begin{tabular}{@{}>{\bfseries\RaggedRight\arraybackslash}p{4.0cm}p{\dimexpr\linewidth-4.0cm-2\tabcolsep\relax}@{}}
    \arrayrulecolor{lightgray}
    Identificador: & \texttt{\l__us_id_tl} \\
    \hline
    Como: & {\l__us_role_tl} \\
    \hline
    Quiero: & {\l__us_goal_tl} \\
    \hline
    Requerimientos: & {\l__us_requirements_tl} \\
    \hline
    Puntos~de~historia: & {\l__us_points_tl} \\
    \hline
    Criterios~de~aceptación: & {\l__us_acceptance_tl} \\
    \hline
    Tareas: & {\l__us_tasks_tl} \\
    \end{tabular}\par
  }
\ExplSyntaxOff


% Referencia a una historia de usuario por su identificador con un enlace clicable
% Ejemplo de uso: \usref{US-XX-YY}
\newcommand{\usref}[1]{\hyperref[us:#1]{\texttt{#1}}}


\begin{document}
% Avoid duplicate page anchors on title/first pages
\hypersetup{pageanchor=false}

\begin{titlepage}
\centering
\includegraphics[width=0.35\textwidth]{logo-uned.jpeg}

\vspace{1.5cm}

{\Large Universidad Nacional de Educación a Distancia\par}
\vspace{0.4cm}
{\large Escuela técnica superior de ingeniería informática\par}

\vspace{1.6cm}

{\large Proyecto de Fin de Grado en Ingeniería Informática\par}

\vspace{1.6cm}

{\LARGE \textbf{Aplicación web de nutrición}\par}

\vspace{2.2cm}

{\large María de la Asunción Gómez Castro\par}

\vfill

{\large Curso 2025-2026\par}
\end{titlepage}
\clearpage
\pagenumbering{arabic}
\setcounter{page}{1}
\hypersetup{pageanchor=true}

\section*{Resumen}
\addcontentsline{toc}{section}{Resumen}
El acceso a información nutricional y la adopción de dietas equilibradas representan un reto cotidiano para muchas personas. Mantener una alimentación variada y saludable, que además considere ingredientes de cercanía y temporada, implica gestionar distintos factores: el valor nutricional de cada alimento, la planificación de menús, y la generación de listas de la compra adecuadas. La dificultad de coordinar todos estos aspectos suele traducirse en hábitos alimentarios poco organizados y, a menudo, menos saludables de lo deseado.

Por otra parte, hoy disponemos de multitud de recursos tecnológicos que pueden facilitar este proceso. Existen bases de datos de alimentos con información detallada, herramientas para digitalizar productos mediante códigos de barras, y una amplia oferta de recetas en línea. Los recientes avances en inteligencia artificial permiten, además, procesar y analizar información de manera eficiente y personalizada.

Este proyecto surge de la necesidad de unificar todas esas herramientas en una única plataforma que simplifique la gestión nutricional y la organización de la cocina diaria. El enfoque consiste en integrar capacidades de importación de recetas, análisis nutricional, planificación de menús y automatización de listas de la compra, haciendo uso de fuentes de datos existentes y de tecnologías como inteligencia artificial, en una aplicación de fácil acceso para usuarios y profesionales de la nutrición.

\newpage

\section*{Abstract}
\addcontentsline{toc}{section}{Abstract}
\begingroup
\selectlanguage{english}
% Disable hyphenation only for the English abstract
\hyphenpenalty=10000\exhyphenpenalty=10000\emergencystretch=2em
Accessing nutritional information and adopting balanced diets is an everyday challenge for many people. Maintaining a varied and healthy diet, while also considering local and seasonal ingredients, involves managing various factors: the nutritional value of each food, menu planning, and generating appropriate shopping lists. The complexity of coordinating all these aspects often results in poorly organized eating habits that are less healthy than intended.

On the other hand, we now have a multitude of technological resources that can help in this process. There are food databases with detailed information, tools for digitizing products using barcode readers, and a wide range of recipes available online. Recent advances in artificial intelligence also enable efficient and personalized processing and analysis of information.

This project was created out of the need to unify all these tools in a single platform that simplifies nutritional management and daily kitchen organization. The approach consists of integrating recipe import capabilities, nutritional analysis, menu planning, and shopping list automation—leveraging existing data sources and technologies like artificial intelligence—into an application that is easily accessible for both individual users and nutrition professionals.
\par
\endgroup
\selectlanguage{spanish}

\newpage
\renewcommand{\contentsname}{Índice}
\addcontentsline{toc}{section}{Índice}
\tableofcontents

\newpage
\section*{Introducción}
\addcontentsline{toc}{section}{Introducción}

Este capítulo analiza el estado del arte en sistemas de nutrición y gestión de recetas, identificando las principales tendencias, aplicaciones y enfoques del sector. Se describe cómo las soluciones tecnológicas actuales abordan problemas relacionados con el control nutricional, la planificación de menús y la gestión de recetas, y se evalúan sus fortalezas y limitaciones. Finalmente, se justifica la necesidad de una nueva solución, Muncher, presentando las funcionalidades que aporta y el hueco que cubre dentro del ecosistema de herramientas disponibles.

\subsection*{Estado del arte}
\addcontentsline{toc}{subsection}{Estado del arte}

En los últimos años, el ecosistema de aplicaciones dedicadas a la nutrición y a la gestión de recetas ha experimentado un fuerte crecimiento, tanto en volumen como en sofisticación. Estas aplicaciones han pasado de ser simples contadores de calorías o libros de recetas digitales a ofrecer herramientas avanzadas de seguimiento dietético, planificación de menús inteligentes, integración con dispositivos móviles y bases de datos alimentarias, y recomendaciones personalizadas. Sin embargo, la mayoría de las soluciones sigue especializándose en aspectos concretos (control nutricional, recetas, planificación, etc.), lo que conduce a una fragmentación de funcionalidades y a la falta de plataformas realmente integradas.

\subsubsection*{Principales actores}
\addcontentsline{toc}{subsubsection}{Principales actores}

A continuación, se resumen las aplicaciones más relevantes en el panorama actual, indicando sus principales características y enlaces a sus sitios oficiales:

\begin{itemize}
  \item \textbf{MyFitnessPal} (\url{https://www.myfitnesspal.com/})\\
  Un referente en conteo de calorías y análisis nutricional. Permite registrar alimentos, crear recetas, fijar objetivos y visualizar progresos a través de gráficos y estadísticas. Incluye una gran base de datos de productos y recetas.

  \item \textbf{Yazio} (\url{https://www.yazio.com/})\\
  Está orientada al seguimiento nutricional y de objetivos de peso, con funcionalidades para el registro de comidas, cálculo de calorías y planificador de comidas.

  \item \textbf{FatSecret} (\url{https://www.fatsecret.es/})\\
  Descripción breve de la aplicación, principales características y público objetivo.

  \item \textbf{Cronometer} (\url{https://cronometer.com/})\\
  Destaca por el seguimiento detallado de micronutrientes, perfecto para dietas médicas o deportivas. Permite personalizar al detalle los objetivos nutricionales.

  \item \textbf{Lifesum} (\url{https://lifesum.com/})\\
  Ofrece seguimiento de hábitos de vida saludables y recomendaciones nutricionales personalizadas. Integra recetas y rutinas.

  \item \textbf{Paprika Recipe Manager} (\url{https://www.paprikaapp.com/})\\
  Aplicación eminentemente gastronómica, con gestión de recetas, planificación de menús y listas de la compra, pero sin funcionalidades nutricionales avanzadas.

  \item \textbf{Mealime} (\url{https://www.mealime.com/})\\
  Planificador de menús semanal y generador de listas de la compra con recetas seleccionadas según preferencias de usuario.

  \item \textbf{Eat This Much} (\url{https://www.eatthismuch.com/})\\
  Planificador automático de menús y recetas alineados con objetivos nutricionales detallados.

  \item \textbf{Tasty} (\url{https://tasty.co/})\\
  Popular por su enfoque visual y social, centrándose en la experiencia culinaria y la inspiración mediante vídeos, pero con poco desarrollo en aspectos nutricionales.
\end{itemize}


\subsubsection*{Enfoques actuales}
\addcontentsline{toc}{subsubsection}{Enfoques actuales}

El conjunto de soluciones disponibles se puede clasificar según tres grandes enfoques, que en ocasiones se combinan parcialmente en una misma aplicación:


\paragraph*{Aplicaciones orientadas a control nutricional}
Estas aplicaciones permiten al usuario registrar lo que come y monitorizar su ingesta de calorías, macronutrientes y micronutrientes. Utilizan grandes bases de datos alimentarias y suelen incluir sugerencias dietéticas, historial de consumo y gráficos de progreso. Ejemplos:

\begin{itemize}
  \item MyFitnessPal
  \item Yazio
  \item FatSecret
  \item Cronometer
  \item Lifesum
\end{itemize}

\paragraph*{Aplicaciones de gestión de recetas y menús}
Estas plataformas se centran en facilitar la organización, edición, planificación y búsqueda de recetas personales o de una comunidad. Permiten al usuario importar, almacenar y modificar recetas, y suelen ofrecer calendarios de menús y listas de la compra automatizadas. Ejemplos:

\begin{itemize}
  \item Paprika
  \item Mealime
  \item Tasty
\end{itemize}

\paragraph*{Aplicaciones de gestión de recetas y menús}
Algunas aplicaciones combinan análisis nutricional y planificación automática, generando menús adaptados a los objetivos del usuario, preferencias alimentarias y restricciones, incluso sugiriendo listas de la compra. Ejemplo destacado: Eat This Much.

\subsubsection*{Proceso para introducir nuevas recetas}
\addcontentsline{toc}{subsubsection}{Proceso para introducir nuevas recetas}

\paragraph*{Entrada manual}
La mayoría de aplicaciones permiten al usuario crear sus propias recetas, pero normalmente exige introducir los ingredientes uno a uno, seleccionando de una base de datos interna y definiendo cantidades y pasos manualmente. Este proceso puede resultar laborioso, especialmente para recetas complejas o largas.

\paragraph*{Importación automática desde la web}
Algunas soluciones (como Paprika y Yummly) permiten importar recetas directamente desde páginas web reconocidas, analizando automáticamente los ingredientes y las instrucciones siempre que sean compatibles con su sistema de reconocimiento.

\paragraph*{Uso de tecnología de IA / procesado de texto}
Si bien hay intentos iniciales de emplear inteligencia artificial para interpretar recetas escritas en texto libre y extraer la información de manera automática, estas capacidades todavía son limitadas. No suele haber integración completa de este tipo de tecnología en la mayoría de soluciones comerciales hasta la fecha.

\paragraph*{Lectores de código de barras}
Apps orientadas a la nutrición (MyFitnessPal, Yazio, FatSecret) permiten registrar productos escaneando el código de barras, facilitando la introducción de ingredientes de productos envasados, aunque dicho escaneo rara vez ayuda a la hora de crear recetas completas.

\subsubsection*{Objetivos de las aplicaciones}
\addcontentsline{toc}{subsubsection}{Objetivos de las aplicaciones}

Además de sus funcionalidades principales, las aplicaciones actuales del ámbito de la nutrición y gestión de recetas suelen estar orientadas a uno o varios objetivos específicos:

\paragraph*{Pérdida de peso}
Muchas de estas herramientas, especialmente las de seguimiento nutricional, están fundamentalmente orientadas a ayudar en el control y la reducción de peso. Proporcionan contadores de calorías, seguimiento del déficit calórico, metas semanales y registros visuales del progreso. Ejemplos:

\begin{itemize}
  \item MyFitnessPal
  \item Yazio
  \item FatSecret
  \item Lifesum
\end{itemize}

\paragraph*{Deporte y objetivos de rendimiento}
Algunas aplicaciones incluyen planes y sugerencias personalizadas para quienes persiguen objetivos deportivos o de rendimiento físico, integrando el seguimiento de macronutrientes (proteínas, carbohidratos, grasas) e incluso el registro de actividad física, a veces conectándose con otros dispositivos o aplicaciones de deporte. Ejemplos:

\begin{itemize}
  \item Cronometer
  \item Lifesum
\end{itemize}

\paragraph*{Salud general y hábitos sostenibles}
Muchas plataformas apuntan también a la mejora de la salud global, afrontando metas como comer más equilibradamente, controlar la ingesta de ciertos nutrientes (por ejemplo, fibra, sodio, azúcares), o cumplir requisitos dietéticos médicos o alergias. Disponen de recordatorios, ajustes de objetivos personalizados e informes de tendencias de salud. Ejemplos:

\begin{itemize}
  \item MyFitnessPal
  \item Yazio
  \item FatSecret
  \item Lifesum
\end{itemize}

\paragraph*{Consejos de estilo de vida y motivación}
Algunas aplicaciones incluyen dentro de su interfaz recomendaciones adicionales que van más allá del plato: consejos de hidratación, actividad física, sueño, mindfulness, y notificaciones motivacionales. En este sentido, buscan fomentar hábitos saludables a través de sugerencias contextualizadas y rutinas complementarias a la alimentación. Ejemplos:

\begin{itemize}
  \item MyFitnessPal
  \item Yazio
  \item FatSecret
  \item Lifesum
\end{itemize}

\paragraph*{Planificación de menús y cocina cotidiana}
Las plataformas orientadas a la gestión de recetas y planificación de menús están más centradas en la organización y variedad de las comidas diarias. Ponen el foco en facilitar la tarea de decidir qué cocinar y cómo adaptar recetas a las preferencias personales, aunque su capacidad de personalización nutricional o integración con consejos de salud suele ser menor que en las apps dirigidas a control nutricional o fitness. Ejemplos:

\begin{itemize}
  \item Paprika
  \item Mealime
  \item Eat This Much
\end{itemize}

\subsection*{Motivación}
\addcontentsline{toc}{subsection}{Motivación}

Pese al avance de las soluciones tecnológicas actuales, persisten limitaciones importantes derivadas de la especialización de cada plataforma. Los usuarios que desean combinar en una sola herramienta la gestión de recetas, el control nutricional detallado, la planificación de menús balanceados y la generación de listas de la compra inteligentes encuentran barreras en forma de procesos laboriosos, duplicidad de registros y falta de integración entre nutrición, cocina y estilo de vida.

\textbf{Muncher} surge como respuesta a este hueco, proponiendo la integración de funcionalidades avanzadas en una interfaz sencilla y unificada. La aplicación pretende facilitar tanto la entrada automática de recetas (mediante importadores inteligentes y tecnologías de IA), como el análisis nutricional preciso y la planificación eficiente del menú semanal, incluyendo sugerencias basadas en preferencias, requisitos nutricionales, estacionalidad y proximidad de ingredientes. El objetivo es que tanto usuarios individuales como nutricionistas puedan disponer de una herramienta realmente útil, flexible y adaptada a la vida cotidiana actual.

\newpage
\section*{Objetivos}
\addcontentsline{toc}{section}{Objetivos}

Esta sección establece los objetivos que guían tanto la elaboración de este documento como el desarrollo de la aplicación \textbf{Muncher}. Se estructura en dos partes diferenciadas pero complementarias.

En primer lugar, se definen los \textbf{objetivos del documento}, que especifican el propósito y alcance de esta memoria técnica, describiendo qué aspectos del proyecto se documentan y con qué nivel de detalle. Estos objetivos orientan la estructura y contenido del presente trabajo, asegurando que se cubran todas las fases relevantes del desarrollo del sistema.

En segundo lugar, se presentan los \textbf{objetivos de la aplicación}, divididos en objetivos funcionales y no funcionales. Los objetivos funcionales describen las capacidades y características que debe ofrecer \textbf{Muncher} a sus usuarios, es decir, qué problemas específicos debe resolver y qué funcionalidades debe proporcionar. Los objetivos no funcionales establecen los criterios de calidad, rendimiento, usabilidad y otros aspectos técnicos que el sistema debe cumplir para garantizar una experiencia satisfactoria y un funcionamiento robusto.

Juntos, estos objetivos proporcionan una hoja de ruta clara tanto para la documentación del proyecto como para su implementación técnica, estableciendo las metas que se pretenden alcanzar y los criterios que permitirán evaluar el éxito del desarrollo.

\subsection*{Objetivos del documento}
\addcontentsline{toc}{subsection}{Objetivos del documento}

El presente documento tiene como finalidad proporcionar una documentación completa y estructurada del proyecto \textbf{Muncher}, abarcando desde la conceptualización inicial hasta la implementación y validación del sistema. Los objetivos específicos que guían la elaboración de esta memoria son los siguientes:

\paragraph*{Establecer una planificación sistemática del proyecto}
mediante la definición clara de los requerimientos funcionales y no funcionales que debe cumplir la aplicación. Esta planificación incluye la clasificación de dichos requerimientos por orden de prioridad, lo que permitirá guiar el desarrollo de manera eficiente y asegurar que las funcionalidades más críticas se implementen en primer lugar.

\paragraph*{Definir los casos de uso más generales}
que representan las interacciones principales entre los usuarios y el sistema. A partir de estos casos de uso, se desarrollará una lista detallada de historias de usuario que permita completar cada funcionalidad de manera incremental y orientada a las necesidades reales de los usuarios finales.

\paragraph*{Documentar la arquitectura del sistema}
de forma justificada, explicando las decisiones de diseño adoptadas y su impacto en el rendimiento, escalabilidad y mantenibilidad de la aplicación. Esta documentación incluirá también la selección y justificación de las herramientas tecnológicas utilizadas en cada capa del sistema.

\paragraph*{Describir el modelo de datos}
necesario para soportar todas las funcionalidades previstas, detallando las entidades, relaciones y restricciones que permitan una implementación robusta y eficiente de la base de datos.

\paragraph*{Explicar las pruebas del sistema}
requeridas para garantizar la funcionalidad, fiabilidad y calidad de la aplicación, incluyendo estrategias de testing que validen tanto los aspectos funcionales como no funcionales del sistema.

\paragraph*{Documentar el proceso de despliegue del sistema en la nube}
explicando las decisiones tomadas respecto a la infraestructura, servicios utilizados, configuración del entorno de producción y estrategias de monitorización. Esta documentación incluirá la justificación de las opciones seleccionadas en términos de escalabilidad, coste y mantenimiento.

\paragraph*{Presentar las conclusiones del proyecto}
evaluando el grado de cumplimiento de los objetivos planteados, analizando las lecciones aprendidas durante el desarrollo y identificando posibles mejoras o líneas de trabajo futuro que puedan enriquecer la aplicación.

Este enfoque metodológico asegura que el documento sirva como guía técnica completa para el desarrollo, implementación y mantenimiento futuro de \textbf{Muncher}.

\newpage
\section*{Requerimientos}
\addcontentsline{toc}{section}{Requerimientos}

Los requerimientos de software constituyen la base fundamental para el desarrollo exitoso de cualquier sistema informático. Según Sommerville \cite{Sommerville2010}, los requerimientos describen los servicios que debe proporcionar el sistema y las restricciones bajo las cuales debe operar, reflejando las necesidades de los clientes.

La ingeniería de requerimientos es un proceso crítico que, según el IEEE Standard 830-1998 \cite{IEEE830-1998}, debe resultar en especificaciones que sean correctas, no ambiguas, completas, consistentes, ordenadas por importancia, verificables, modificables y trazables. Como indica el Project Management Institute \cite{PMI2021}, una gestión ineficaz de los requerimientos puede llevar a retrabajos, expansión del alcance, insatisfacción del cliente, sobrecostes, retrasos en el cronograma y, en general, al fracaso del proyecto.

Los requerimientos se clasifican tradicionalmente en dos categorías principales: requerimientos funcionales, que describen el comportamiento observable del sistema bajo diversas circunstancias, y requerimientos no funcionales, que especifican características importantes del sistema como disponibilidad, usabilidad, seguridad o rendimiento (Wiegers y Beatty \cite{Wiegers2013}).

En esta sección se presentan los requerimientos específicos para \textbf{Muncher}, clasificados según su naturaleza y priorizados de acuerdo con su importancia para el éxito del proyecto. Esta especificación detallada servirá como guía para el diseño, implementación y validación del sistema.

\subsection*{Priorización de requerimientos}
\addcontentsline{toc}{subsection}{Priorización de requerimientos}

La priorización de requerimientos es un proceso fundamental en el desarrollo de software que consiste en ordenar los requerimientos según su importancia relativa para el éxito del proyecto. Como señalan Karlsson y Ryan \cite{Karlsson1997}, "un conocimiento claro e inequívoco sobre las prioridades de los requerimientos ayuda a enfocar el proceso de desarrollo y gestionar los proyectos de manera más efectiva y eficiente".

La importancia de una priorización efectiva radica en que no todos los requerimientos tienen el mismo impacto en el valor del producto final. Una priorización adecuada permite tomar decisiones informadas sobre qué características implementar primero, facilita la gestión de cambios durante el desarrollo y ayuda a establecer una hoja de ruta clara para las entregas del proyecto.

\subsubsection*{Método MoSCoW}
\addcontentsline{toc}{subsubsection}{Método MoSCoW}

Para la priorización de los requerimientos de Muncher se utilizará el método MoSCoW, una técnica de priorización desarrollada como parte del framework DSDM (Dynamic Systems Development Method) \cite{DSDM2014} que ayuda a comprender y gestionar prioridades de manera efectiva.

Según el DSDM Consortium, en proyectos donde el tiempo es un factor limitante, es vital comprender la importancia relativa del trabajo a realizar para mantener el progreso y cumplir con los plazos establecidos. El método MoSCoW clasifica los requerimientos en cuatro categorías claramente definidas:

\begin{itemize}
    \item \textbf{Must Have (M):} Los requerimientos Must Have constituyen el Subconjunto Mínimo Utilizable (MUST) que el proyecto garantiza entregar, siendo absolutamente esenciales para que la solución sea viable, legal o segura. La prueba decisiva para identificarlos es preguntarse: "¿Qué sucede si este requerimiento no se cumple?"; si la respuesta implica cancelar el proyecto porque no tendría sentido implementar la solución sin él, entonces es un Must Have, mientras que si existe alguna alternativa (aunque sea manual y laboriosa), pertenece a categorías de menor prioridad.
    \item \textbf{Should Have (S):} Los requerimientos Should Have se definen como importantes pero no vitales: puede ser doloroso omitirlos, pero la solución sigue siendo viable. Pueden requerir algún tipo de solución alternativa como gestión de expectativas, cierta ineficiencia, una solución existente o trabajo manual, que puede ser temporal. La diferencia clave entre un Should Have y un Could Have radica en el grado de dolor causado por no cumplir el requerimiento, medido en términos de valor de negocio o número de personas afectadas.
    \item \textbf{Could Have (C):} Los requerimientos Could Have se definen como deseables pero menos importantes, con menor impacto si se omiten comparado con los Should Have. Estos requerimientos proporcionan la principal reserva de contingencia del proyecto, ya que solo se entregarían en su totalidad en el mejor escenario posible. Cuando surge un problema y la fecha límite está en riesgo, los Could Have constituyen la primera opción de elementos a eliminar del marco temporal actual.
    \item \textbf{Won't Have This Time (W):} Los requerimientos Won't Have This Time son aquellos que el equipo del proyecto ha acordado que no serán entregados como parte de este marco temporal. Se registran en la Lista de Requerimientos Priorizados donde ayudan a clarificar el alcance del proyecto, evitando que sean reintroducidos informalmente en una fecha posterior. Esto ayuda a gestionar las expectativas de que algunos requerimientos simplemente no llegarán a la Solución Desplegada, al menos no en esta ocasión.
\end{itemize}



\subsection*{Requerimientos funcionales}
\addcontentsline{toc}{subsection}{Requerimientos funcionales}

Los requerimientos funcionales especifican las funcionalidades concretas que debe proporcionar el sistema para satisfacer las necesidades identificadas de los usuarios. Estos requerimientos describen las capacidades específicas, comportamientos y servicios que \textbf{Muncher} debe ofrecer, estableciendo qué debe hacer el sistema sin entrar en detalles sobre cómo debe hacerlo.

La especificación detallada de los requerimientos funcionales es fundamental para garantizar que el sistema desarrollado cumpla con las expectativas de los usuarios y proporcione valor real en la resolución de los problemas identificados. Para facilitar su gestión y trazabilidad a lo largo del proyecto, los requerimientos se han organizado en categorías temáticas coherentes con los objetivos funcionales previamente establecidos:

\begin{itemize}
    \item \textbf{Gestión de usuarios (GU):} Requerimientos relacionados con la gestión de usuarios y su acceso al sistema.
    \item \textbf{Introducción de recetas (IR):} Requerimientos relacionados con la introducción de recetas en el sistema.
    \item \textbf{Planificación de menús (PM):} Requerimientos relacionados con la planificación de menús.
    \item \textbf{Consulta de valores nutricionales (VN):} Requerimientos relacionados con la consulta de valores nutricionales.
    \item \textbf{Gestión de listas de la compra (LC):} Requerimientos relacionados con la gestión de listas de la compra.
\end{itemize}

Cada requerimiento funcional ha sido identificado con un código único siguiendo el formato \textbf{RF-XX-YY}, donde \textbf{RF} indica "Requerimiento Funcional", \textbf{XX} indica el tipo de requerimiento funcional y \textbf{YY} corresponde a un número secuencial. Esta codificación permite la trazabilidad completa desde la especificación inicial hasta la implementación, pruebas y validación final, facilitando la gestión de cambios y el seguimiento del progreso del desarrollo.



\subsubsection*{Gestión de usuarios}
\addcontentsline{toc}{subsubsection}{Gestión de usuarios}

\reqtitle{RF-GU-01}{Registro de usuarios}
\req{RF-GU-01}{El sistema permitirá a los usuarios registrarse en la aplicación.}{Must Have}

\reqtitle{RF-GU-02}{Inicio de sesión}
\req{RF-GU-02}{El sistema permitirá a los usuarios iniciar sesión en la aplicación.}{Must Have}

\reqtitle{RF-GU-03}{Cierre de sesión}
\req{RF-GU-03}{El sistema permitirá a los usuarios cerrar sesión en la aplicación.}{Must Have}

\reqtitle{RF-GU-04}{Administración de usuarios}
\req{RF-GU-04}{El sistema permitirá a los usuarios administradores gestionar los usuarios de la aplicación.}{Should Have}

\reqtitle{RF-GU-05}{Eliminación de cuentas de usuario}
\req{RF-GU-05}{El sistema permitirá a los usuarios administradores eliminar cuentas de usuario.}{Won't Have This Time}

\subsubsection*{Introducción de recetas}
\addcontentsline{toc}{subsubsection}{Introducción de recetas}

\reqtitle{RF-IR-01}{Introducción manual de recetas}
\req{RF-IR-01}{El sistema permitirá a los usuarios introducir recetas manualmente.}{Must Have}

\reqtitle{RF-IR-02}{Análisis automático de recetas}
\req{RF-IR-02}{El sistema analizará automáticamente recetas de fuentes externas y las incorporará al inventario del usuario.}{Should Have}

\reqtitle{RF-IR-03}{Clasificación de recetas}
\req{RF-IR-03}{El sistema permitirá a los usuarios clasificar las recetas en categorías.}{Could Have}

\reqtitle{RF-IR-04}{Búsqueda de recetas}
\req{RF-IR-04}{El sistema permitirá a los usuarios buscar recetas por nombre, ingredientes o categorías.}{Must Have}

\reqtitle{RF-IR-05}{Edición de recetas}
\req{RF-IR-05}{El sistema permitirá a los usuarios editar las recetas.}{Must Have}

\reqtitle{RF-IR-06}{Eliminación de recetas}
\req{RF-IR-06}{El sistema permitirá a los usuarios eliminar recetas.}{Must Have}

\reqtitle{RF-IR-07}{Visualización de recetas}
\req{RF-IR-07}{El sistema permitirá a los usuarios visualizar las recetas.}{Must Have}

\reqtitle{RF-IR-08}{Creación de variantes de recetas}
\req{RF-IR-08}{El sistema permitirá a los usuarios crear variantes de recetas existentes para adaptarlas a diferentes escenarios.}{Could Have}

\reqtitle{RF-IR-09}{Enriquecimiento de recetas}
\req{RF-IR-09}{El sistema enriquecerá las recetas con información adicional como las dietas en las que se puede incluir o las precauciones de los ingredientes.}{Could Have}

\subsubsection*{Planificación de menús}
\addcontentsline{toc}{subsubsection}{Planificación de menús}

\reqtitle{RF-PM-01}{Planificación manual de menú}
\req{RF-PM-01}{El sistema permitirá a los usuarios planificar menús añadiendo recetas manualmente.}{Must Have}

\reqtitle{RF-PM-02}{Planificación automática de menú}
\req{RF-PM-02}{El sistema permitirá a los usuarios generar un menú planificado automáticamente en base a una serie de parámetros: ingredientes disponibles, restricciones alimentarias, objetivos, restricciones de tiempo, etc.}{Could Have}

\reqtitle{RF-PM-03}{Visualización de menús}
\req{RF-PM-03}{El sistema permitirá a los usuarios visualizar los menús planificados.}{Must Have}

\reqtitle{RF-PM-04}{Edición de menús}
\req{RF-PM-04}{El sistema permitirá a los usuarios editar los menús planificados.}{Must Have}

\reqtitle{RF-PM-05}{Eliminación de menús}
\req{RF-PM-05}{El sistema permitirá a los usuarios eliminar los menús planificados.}{Must Have}

\subsubsection*{Consulta de valores nutricionales}
\addcontentsline{toc}{subsubsection}{Consulta de valores nutricionales}

\reqtitle{RF-VN-01}{Consulta de valores nutricionales de una receta}
\req{RF-VN-01}{El sistema permitirá a los usuarios consultar los valores nutricionales de una receta.}{Must Have}

\reqtitle{RF-VN-02}{Consulta de valores nutricionales de un periodo de tiempo}
\req{RF-VN-02}{El sistema permitirá a los usuarios consultar los valores nutricionales de un periodo de tiempo, como un día o una semana.}{Must Have}

\reqtitle{RF-VN-03}{Establecimiento manual de objetivos nutricionales}
\req{RF-VN-03}{El sistema permitirá a los usuarios establecer objetivos nutricionales manualmente.}{Should Have}

\reqtitle{RF-VN-04}{Establecimiento de objetivos nutricionales automáticos}
\req{RF-VN-04}{El sistema permitirá a los usuarios establecer objetivos nutricionales automáticamente en base a su peso, altura, edad, objetivos, etc.}{Could Have}

\reqtitle{RF-VN-05}{Visualización de progreso en la consecución de objetivos nutricionales}
\req{RF-VN-05}{El sistema permitirá a los usuarios visualizar el progreso en la consecución de sus objetivos nutricionales.}{Should Have}

\subsubsection*{Gestión de listas de la compra}
\addcontentsline{toc}{subsubsection}{Gestión de listas de la compra}

\reqtitle{RF-LC-01}{Creación manual de listas de la compra}
\req{RF-LC-01}{El sistema permitirá a los usuarios crear listas de la compra manualmente.}{Must Have}

\reqtitle{RF-LC-02}{Creación automática de listas de la compra}
\req{RF-LC-02}{El sistema permitirá a los usuarios crear listas de la compra automáticamente en base a los ingredientes necesarios para las recetas planificadas.}{Must Have}

\reqtitle{RF-LC-03}{Visualización de listas de la compra}
\req{RF-LC-03}{El sistema permitirá a los usuarios visualizar las listas de la compra.}{Must Have}

\reqtitle{RF-LC-04}{Edición de listas de la compra}
\req{RF-LC-04}{El sistema permitirá a los usuarios editar las listas de la compra.}{Must Have}

\reqtitle{RF-LC-05}{Eliminación de listas de la compra}
\req{RF-LC-05}{El sistema permitirá a los usuarios eliminar las listas de la compra.}{Must Have}

\reqtitle{RF-LC-06}{Progreso de listas de la compra}
\req{RF-LC-06}{El sistema permitirá a los usuarios marcar los ingredientes ya disponibles o comprados en la lista de la compra a medida que se van comprando.}{Should Have}

\subsection*{Requerimientos no funcionales}
\addcontentsline{toc}{subsection}{Requerimientos no funcionales}

Los requerimientos no funcionales especifican las cualidades y restricciones que debe cumplir el sistema en términos de rendimiento, usabilidad, seguridad y otros atributos de calidad. A diferencia de los requerimientos funcionales, que describen qué debe hacer el sistema, los requerimientos no funcionales establecen cómo debe comportarse y bajo qué condiciones debe operar para proporcionar una experiencia satisfactoria a los usuarios.

Estos requerimientos son fundamentales para el éxito del proyecto, ya que determinan aspectos críticos como la facilidad de uso, la confiabilidad, la seguridad y el rendimiento del sistema. Su correcta especificación y cumplimiento puede marcar la diferencia entre una aplicación técnicamente funcional y una verdaderamente exitosa en el mercado.

Para facilitar su organización y trazabilidad, los requerimientos no funcionales se han estructurado en las siguientes categorías, alineadas con los objetivos no funcionales previamente establecidos:

\begin{itemize}
    \item \textbf{Compatibilidad (CO):} Requerimientos relacionados con la utilización en diferentes dispositivos.
    \item \textbf{Inteligencia Artificial (IA):} Requerimientos sobre la integración de herramientas de IA.
    \item \textbf{Calidad (CA):} Requerimientos de automatización de controles de calidad.
    \item \textbf{Despliegue (DE):} Requerimientos de automatización del despliegue.
    \item \textbf{Monitorización (MO):} Requerimientos de seguimiento del estado del sistema.
    \item \textbf{Seguridad (SE):} Requerimientos de protección y control de acceso.
    \item \textbf{Rendimiento (RE):} Requerimientos de velocidad y eficiencia.
    \item \textbf{Usabilidad (US):} Requerimientos de facilidad de uso y experiencia de usuario.
\end{itemize}

Cada requerimiento no funcional ha sido identificado con un código único siguiendo el formato \textbf{RNF-XX-YY}, donde RNF indica "Requerimiento No Funcional", XX corresponde al código de categoría y YY a un número secuencial dentro de esa categoría.

\subsubsection*{Compatibilidad}
\addcontentsline{toc}{subsubsection}{Compatibilidad}

\reqtitle{RNF-CO-01}{Tecnologías web}
\req{RNF-CO-01}{El sistema utilizará tecnologías web estándar, que facilitarán su acceso desde cualquier dispositivo con navegación web.}{Must Have}

\reqtitle{RNF-CO-02}{Responsividad}
\req{RNF-CO-02}{El sistema será responsivo, adaptándose a diferentes tamaños de pantalla y resoluciones.}{Must Have}

\subsubsection*{Inteligencia Artificial}
\addcontentsline{toc}{subsubsection}{Inteligencia Artificial}

\reqtitle{RNF-IA-01}{Introducción automática de datos}
\req{RNF-IA-01}{El sistema utilizará tecnologías de IA para introducir datos de forma automática a partir de imágenes o texto en lenguaje natural.}{Must Have}

\reqtitle{RNF-IA-02}{Generación automática de sugerencias}
\req{RNF-IA-02}{El sistema utilizará tecnologías de IA para generar sugerencias de recetas, menús, ingredientes, etc.}{Could Have}

\subsubsection*{Calidad}
\addcontentsline{toc}{subsubsection}{Calidad}

\reqtitle{RNF-CA-01}{Integración continua}
\req{RNF-CA-01}{El sistema tendrá integración continua, para validar la correcta funcionalidad del mismo.}{Must Have}

\reqtitle{RNF-CA-02}{Control de calidad de código}
\req{RNF-CA-02}{El sistema tendrá un control de calidad de formato y buenas prácticas de código fuente, que se usará para validar los cambios antes de que se integren en el repositorio.}{Must Have}

\reqtitle{RNF-CA-03}{Control de calidad de funcionalidad}
\req{RNF-CA-03}{El sistema tendrá un control de calidad de funcionalidad, que se usará para validar que el sistema cumple con los requerimientos funcionales.}{Must Have}

\reqtitle{RNF-CA-04}{Análisis de gramática y ortografía}
\req{RNF-CA-04}{El sistema tendrá un análisis de gramática y ortografía, que se usará para validar que el texto introducido es correcto.}{Should Have}

\subsubsection*{Despliegue}
\addcontentsline{toc}{subsubsection}{Despliegue}

\reqtitle{RNF-DE-01}{Despliegue en la nube}
\req{RNF-DE-01}{El sistema se desplegará en la nube, para poder acceder a él desde cualquier dispositivo con acceso a internet.}{Must Have}

\reqtitle{RNF-DE-02}{Despliegue continuo}
\req{RNF-DE-02}{Cada nuevo cambio en el código fuente se desplegará a la nube de forma automática.}{Must Have}

\reqtitle{RNF-DE-03}{Separación de entornos}
\req{RNF-DE-03}{Existirán dos entornos de nube diferentes: uno para desarrollo y otro para producción.}{Could Have}

\reqtitle{RNF-DE-04}{Entorno local}
\req{RNF-DE-04}{El proyecto incluirá mecanismos e instrucciones para levantar localmente un entorno idéntico a la nube que el desarrollador pueda usar para probar el sistema antes de enviar sus cambios en el código.}{Must Have}

\subsubsection*{Monitorización}
\addcontentsline{toc}{subsubsection}{Monitorización}

\reqtitle{RNF-MO-01}{Monitorización de errores}
\req{RNF-MO-01}{El sistema tendrá un sistema de monitorización de errores, para poder detectar y solucionar los errores que se produzcan.}{Must Have}

\reqtitle{RNF-MO-02}{Alertas ante errores}
\req{RNF-MO-02}{El sistema tendrá un sistema de alertas ante errores, para poder notificar a los desarrolladores cuando se produzcan errores.}{Should Have}

\reqtitle{RNF-MO-03}{Monitorización de rendimiento}
\req{RNF-MO-03}{El sistema tendrá un sistema de monitorización de rendimiento, para poder detectar y solucionar los problemas de rendimiento que se produzcan. Incluirá trazas que informen, para cada operación, del tiempo invertido en ella.}{Could Have}

\reqtitle{RNF-MO-04}{Monitorización de tráfico}
\req{RNF-MO-04}{El sistema tendrá un sistema de monitorización de tráfico, para conocer en todo momento el número de usuarios que están usando el sistema.}{Could Have}

\subsubsection*{Seguridad}
\addcontentsline{toc}{subsubsection}{Seguridad}

\reqtitle{RNF-SE-01}{JSON Web Tokens}
\req{RNF-SE-01}{El sistema utilizará JSON web tokens para verficar la identidad de los usuarios.}{Must Have}

\reqtitle{RNF-SE-02}{Encriptación}
\req{RNF-SE-02}{El sistema encriptará las credenciales del usuario antes de almacenarlas entre sesiones.}{Must Have}

\reqtitle{RNF-SE-03}{Almacenamiento seguro de secretos}
\req{RNF-SE-03}{Los datos sensibles del sistema, como las claves de acceso a la nube o la contraseña de la base de datos, se almacenarán de forma segura.}{Must Have}

\reqtitle{RNF-SE-04}{Monitorización de vulnerabilidades}
\req{RNF-SE-04}{El sistema tendrá un sistema de monitorización de vulnerabilidades, para poder detectar y solucionar las vulnerabilidades que se produzcan.}{Should Have}

\subsubsection*{Rendimiento}
\addcontentsline{toc}{subsubsection}{Rendimiento}

\reqtitle{RNF-RE-01}{Uso de SPA}
\req{RNF-RE-01}{La interfaz de usuario se implementará como una SPA (Single Page Application), para asegurar que el feedback a las interacciones es inmediato.}{Must Have}

\reqtitle{RNF-RE-02}{Caché de datos en el front-end}
\req{RNF-RE-02}{El front-end del sistema tendrá una caché interna de los datos recibidos del back-end, para ahorrar tiempo en llamadas sucesivas.}{Must Have}

\reqtitle{RNF-RE-03}{Servicio de índice de datos}
\req{RNF-RE-03}{El sistema tendrá un servicio de índice de datos sobre la base de datos, para optimizar el tiempo de consultas.}{Should Have}

\reqtitle{RNF-RE-04}{Red de distribución de contenido}
\req{RNF-RE-04}{Los ficheros estáticos de la aplicación, como imágenes o vídeos, se distribuirán mediante una CDN (Content Delivery Network), para reducir la latencia de las peticiones.}{Could Have}

\subsubsection*{Usabilidad}
\addcontentsline{toc}{subsubsection}{Usabilidad}

\reqtitle{RNF-US-01}{Diseño minimalista}
\req{RNF-US-01}{El sistema tendrá un diseño minimalista, para que se pueda usar sin necesidad de instrucciones.}{Must Have}

\reqtitle{RNF-US-02}{Validación de formularios}
\req{RNF-US-02}{Los formularios incluirán validación en tiempo real con mensajes de error claros y específico.}{Should Have}

\reqtitle{RNF-US-03}{Legibilidad}
\req{RNF-US-03}{Los textos tendrán suficiente contraste de color para garantizar la legibilidad
.}{Must Have}


\newpage
\section*{Bibliografía}
\addcontentsline{toc}{section}{Bibliografía}

\begin{thebibliography}{9}
\bibitem{Sommerville2010}
Sommerville, I. (2010).
\textit{Software Engineering} (9th ed., p. 83).
Addison-Wesley.
\url{https://engineering.futureuniversity.com/BOOKS%20FOR%20IT/Software-Engineering-9th-Edition-by-Ian-Sommerville.pdf}
\bibitem{IEEE830-1998}
IEEE Computer Society. (1998).
\textit{IEEE Standard for Software Requirements Specifications} (IEEE Std 830-1998, p. 4).
IEEE Computer Society.
\url{http://www.math.uaa.alaska.edu/~afkjm/cs401/IEEE830.pdf}
\bibitem{PMI2021}
Project Management Institute. (2021)
\textit{A Guide to the Project Management Body of Knowledge (PMBOK Guide)} (7th ed., p. 83).
PMI Publications.
\url{https://ibimone.com/PMBOK%207th%20Edition%20(iBIMOne.com).pdf}
\bibitem{Wiegers2013}
Wiegers, K. E., \& Beatty, J. (2013).
\textit{Software Requirements} (3rd ed., p. 10).
Microsoft Press.
\url{https://olivroqueaprende.com/WDK/Software_Requirements_3rd_Edition.pdf}
\bibitem{Karlsson1997}
Karlsson, A., \& Ryan, C. (1997).
\textit{A Cost–Value Approach for Prioritizing Requirements} (2nd ed., p. 68).
IEEE SOFTWARE.
\url{https://fileadmin.cs.lth.se/cs/Education/ETSF30/Labb/karlsson97.pdf}
\bibitem{DSDM2014}
Dynamic Systems Development Method Consortium. (2014).
\textit{The DSDM Agile Project Framework Handbook.} (c. 10).
DSDM Consortium.
\url{https://www.agilebusiness.org/dsdm-project-framework/moscow-prioririsation.html}
\end{thebibliography}

\newpage
\section*{Anexo I: Historias de usuario}
\phantomsection
\addcontentsline{toc}{section}{Anexo I: Historias de usuario}
\label{sec:anexo-historias-usuario}

Las historias de usuario constituyen una técnica fundamental en el desarrollo ágil de software para capturar requerimientos desde la perspectiva del usuario final. Según Cohn \cite{Cohn2004}, "una historia de usuario describe funcionalidad que será valiosa para un usuario o cliente de un sistema o software", y están compuestas por tres aspectos: una descripción escrita utilizada para planificación y como recordatorio, conversaciones sobre la historia que ayudan a elaborar los detalles, y pruebas que pueden usarse para determinar cuándo una historia está completa1. Estas historias sirven como herramienta para alcanzar tanto requerimientos funcionales como no funcionales, funcionando como pasos individuales e incrementales que van acercando el desarrollo hacia el objetivo final del proyecto.

Cada historia de usuario representa una funcionalidad específica que aporta valor al usuario y puede ser implementada de forma independiente en una iteración del proyecto. La asignación de estas historias a las diferentes iteraciones se realizará en función de las prioridades establecidas a los requerimientos con los que se asocian, asegurando que las funcionalidades más críticas se desarrollen en primer lugar.

Para facilitar la planificación y estimación del esfuerzo de desarrollo, cada historia de usuario será valorada utilizando puntos de historia basados en la sucesión de Fibonacci (1, 2, 3, 5, 8). Como explica Cohn \cite{Cohn2005}, los estudios demuestran que somos mejores estimando cosas que caen dentro de un orden de magnitud, y estas secuencias no lineales funcionan bien porque reflejan la mayor incertidumbre asociada con estimaciones para unidades de trabajo más grandes. Los espacios en la secuencia de Fibonacci se vuelven apropiadamente más grandes a medida que los números aumentan, donde cada número debe pensarse como un "cajón" en el que se vierten elementos del tamaño apropiado. Los puntos no representan tiempo directamente, sino una medida relativa de complejidad, esfuerzo y riesgo asociado a cada historia.

Cada historia de usuario desarrollada estará explícitamente relacionada con uno o varios requerimientos funcionales o no funcionales identificados previamente, estableciendo una trazabilidad clara entre las necesidades documentadas y su implementación práctica. Esta relación garantiza que cada incremento de funcionalidad contribuya de manera medible al cumplimiento de los objetivos del proyecto.

\subsection*{Configuración inicial}
\addcontentsline{toc}{subsection}{Configuración inicial}

\us[
  identificador=US-CI-01,
  titulo={Creación del repositorio de código},
  como={Desarrollador},
  quiero={Controlar las versiones de mi código, compartirlo con otros compañeros y mantener copias de seguridad para mejorar el acceso al mismo},
  requerimientos={\reqref{RNF-CA-01}, \reqref{RNF-DE-02}},
  puntos={1},
  criterios={
    \begin{itemize}
      \item El repositorio de código en la nube existe.
      \item Es accesible por todos los desarrolladores.
      \item Incluye instrucciones de uso.
    \end{itemize}
  },
  tareas={
    \begin{itemize}
      \item Inicialización de Git en un directorio local.
      \item Creación de un repositorio en Github.
      \item Configuración de acceso remoto al repositorio de Github desde la máquina local.
      \item Redacción de las instrucciones de uso en un README.
    \end{itemize}
  }
]

\us[
  identificador=US-CI-02,
  titulo={Creación de una SPA usando Vite y React},
  como={Desarrollador},
  quiero={Ver mi proyecto en funcionamiento para poder verificar que los cambios que realizo tienen el efecto deseado},
  requerimientos={\reqref{RNF-CO-01}, \reqref{RNF-DE-04}, \reqref{RNF-RE-01}},
  puntos={2},
  criterios={
    \begin{itemize}
      \item Existe una aplicación Vite + React.
      \item La aplicación puede ejecutarse localmente.
      \item La aplicación puede visitarse en el navegador.
      \item La aplicación responde de forma automática a los cambios llevados a cabo en el código fuente.
    \end{itemize}
  },
  tareas={
    \begin{itemize}
      \item Creación de un proyecto Vite + React.
      \item Personalización de la configuración.
    \end{itemize}
  }
]

\end{document}